% ============================================================
%  STRUCTURALLY INCORRUPTIBLE TARGETED ACQUISITION IN
%  MASS SPECTROMETRY: A TEMPORAL PROGRAMMING FRAMEWORK
% ============================================================

\documentclass[twocolumn,10pt,a4paper]{article}

\usepackage[utf8]{inputenc}
\usepackage[T1]{fontenc}
\usepackage{lmodern}
\usepackage[a4paper,top=2cm,bottom=2.2cm,left=1.8cm,right=1.8cm,columnsep=0.6cm]{geometry}
\usepackage{amsmath,amssymb,amsthm,mathtools}
\usepackage{bm}
\usepackage{booktabs}
\usepackage{array}
\usepackage{enumitem}
\usepackage{graphicx}
\usepackage{algorithm}
\usepackage{algpseudocode}
\usepackage{xcolor}
\usepackage{microtype}
\usepackage[round,sort&compress]{natbib}
\usepackage[colorlinks=true,linkcolor=blue!60!black,
            citecolor=blue!60!black,urlcolor=blue!60!black]{hyperref}
\usepackage{fancyhdr}

\pagestyle{fancy}
\fancyhf{}
\fancyhead[LE,RO]{\thepage}
\fancyhead[RE]{\small Structurally Incorruptible Targeted Acquisition}
\fancyhead[LO]{\small Sachikonye}
\renewcommand{\headrulewidth}{0.4pt}

% ---- Theorem environments -------------------------------------------
\newtheoremstyle{thmstyle}{\topsep}{\topsep}{\itshape}{}{\bfseries}{.}{.5em}{}
\theoremstyle{thmstyle}
\newtheorem{theorem}{Theorem}[section]
\newtheorem{lemma}[theorem]{Lemma}
\newtheorem{proposition}[theorem]{Proposition}
\newtheorem{corollary}[theorem]{Corollary}
\theoremstyle{definition}
\newtheorem{definition}[theorem]{Definition}
\newtheorem{axiom}{Axiom}
\newtheorem{example}[theorem]{Example}
\theoremstyle{remark}
\newtheorem{remark}[theorem]{Remark}

% ---- Custom commands ------------------------------------------------
\newcommand{\R}{\mathbb{R}}
\newcommand{\N}{\mathbb{N}}
\newcommand{\Z}{\mathbb{Z}}
\newcommand{\kB}{k_{\mathrm{B}}}
\newcommand{\DP}{\Delta P}
\newcommand{\Tref}{T_{\mathrm{ref}}}
\newcommand{\trec}{t_{\mathrm{rec}}}
\newcommand{\Mcyc}{M}
\newcommand{\fref}{f_{\mathrm{ref}}}
\newcommand{\calC}{\mathcal{C}}
\newcommand{\calD}{\mathcal{D}}
\newcommand{\calA}{\mathcal{A}}
\newcommand{\calP}{\mathcal{P}}
\newcommand{\Leb}{\mathrm{Leb}}
\newcommand{\LM}{\mathcal{L}_{\mathcal{M}}}
\newcommand{\AM}{\mathbf{A}_{\mathcal{M}}}
\newcommand{\Mop}{\mathcal{M}}
\newcommand{\tp}{\langle\tau_{p}\rangle}

% =========================================================
\title{\textbf{Structurally Incorruptible Targeted Acquisition\\
in Mass Spectrometry: A Temporal Programming Framework}}

\author{
Kundai Farai Sachikonye\\
Technical University of Munich\\
Chair of Proteomics and Bioanalytics\\
\texttt{kundai.sachikonye@wzw.tum.de}
}
\date{\today}
% =========================================================

\begin{document}
\maketitle

% =========================================================
\begin{abstract}
\noindent
We develop a unified framework for real-time targeted acquisition in
mass spectrometry grounded in two independently motivated structures:
the \emph{bounded phase space formulation} of ion dynamics, which shows
that all four mass analyzer equations are Euler-Lagrange consequences
of a single partition Lagrangian; and \emph{temporal programming}, a
paradigm in which every inter-node signal is a bare pulse and the sole
computational datum is the timing deviation
$\DP(k) = \Tref(k) - \trec(k)$.

We prove that mass spectrometry is intrinsically a timing instrument:
$m/z$ is never measured directly but inferred from a time (TOF) or a
frequency (Orbitrap, FT-ICR). The natural control variable is therefore
$\DP$, not $m/z$. A temporal programming acquisition compiles the target
list into a \emph{cell registry} $\Gamma$ over $\DP$-space and makes
real-time accept/discard decisions by cell membership. Because no payload
is parsed, we prove a \emph{Structural Incorruptibility Theorem}: no
off-target ion can satisfy any calibrated cell condition, and the
framework has zero injection surface by construction.

We derive from first principles the \emph{partition uncertainty relation}
$\Delta M \cdot \tp \geq \hbar$, which governs the tradeoff between
real-time trigger speed and mass resolution. This yields an exact design
equation: for a target cell of width $\delta_{m/z}$, the minimum
acquisition time before the trigger can fire is
$\tau_{\min} = \hbar/\Delta M_{\mathrm{cell}}$, derivable from the
partition Lagrangian alone.

For multi-dimensional acquisition, we prove the
\emph{Composition-Inflation Theorem}:
$T(n,d) = d(d+1)^{n-1}$ distinguishable ion states are accessible to a
system with $d$ timing channels (dimensions) and $n$ timing events. Ion
mobility drift time constitutes a natural second timing axis; a two-channel
temporal programming system resolves $T(n,2) = 2 \cdot 3^{n-1}$ states,
separating isobars that are inaccessible to single-axis timing.

\medskip\noindent\textbf{Keywords:} targeted acquisition, temporal
programming, timing cells, partition Lagrangian, structural incorruptibility,
real-time trigger, partition uncertainty, composition-inflation, ion mobility,
SRM, PRM, DDA
\end{abstract}

\tableofcontents

% =========================================================
\section{Introduction}
\label{sec:intro}
% =========================================================

\subsection{The targeted acquisition problem}

Mass spectrometry has become the principal analytical platform for
quantitative proteomics, metabolomics, and pharmaceutical analysis. The
fundamental measurement task is to identify and quantify molecular ions
by their mass-to-charge ratio $m/z$. Three acquisition strategies
dominate practice.

\emph{Data-dependent acquisition} (DDA) performs a survey scan, ranks
precursor ions by intensity, and selects the most abundant for
fragmentation. The selection is post hoc: the instrument acquires a full
spectrum, interprets it, and acts on the interpretation. This introduces
a parsing step that is both slow (duty cycles of 0.1--1\,s) and
vulnerable to chemical noise overwhelming genuine signals.

\emph{Selected/parallel reaction monitoring} (SRM/PRM) configures the
mass filter to pass only pre-specified $m/z$ windows. This is highly
specific but operates on the mass scale, not the time scale; the filter
is a mechanical or electromagnetic gate that is set per-target and
requires prior knowledge of the exact $m/z$ value.

\emph{Data-independent acquisition} (DIA/SWATH) fragments all ions in a
wide window, acquiring the full fragment spectrum for downstream
disambiguation \citep{gillet2012}. This is comprehensive but produces
highly chimeric spectra requiring complex deconvolution.

All three strategies share a structural assumption: the primary control
variable is $m/z$. We show this assumption is unnecessary and
disadvantageous. $m/z$ is never the direct observable in any analyzer;
it is always reconstructed from a timing measurement. The natural control
variable is the timing deviation $\DP(k)$, and building the acquisition
paradigm around $\DP$ yields a system that is faster, more specific, and
provably immune to a class of failures that $m/z$-based acquisition
cannot prevent.

\subsection{Mass spectrometry is a timing instrument}

This observation, while known informally, has not been made the
foundation of an acquisition paradigm. Consider the four commercial
analyzer types:

\begin{itemize}[nosep]
\item \textbf{Time-of-flight (TOF)}: ions are separated by flight time
  $T \propto \sqrt{m/z}$. The physical observable is $T$.
\item \textbf{Orbitrap}: ions oscillate axially at frequency
  $\omega_z \propto \sqrt{z/m}$. The physical observable is $\omega_z$,
  recovered by Fourier transform of the transient image current.
\item \textbf{FT-ICR}: ions undergo cyclotron motion at
  $\omega_c \propto z/m$. Again the observable is a frequency, recovered
  from a transient.
\item \textbf{Quadrupole mass filter}: ions traverse a stability boundary
  defined by the Mathieu parameters $a_u = 4qU/(m\Omega^2 r_0^2)$ and
  $q_u = 2qV/(m\Omega^2 r_0^2)$. Stability is determined by the
  ratio $\Omega t$, a dimensionless time.
\end{itemize}

In every case, $m/z$ is inferred from a timing or frequency measurement.
The $m/z$ value never appears directly on any sensor. This makes all
mass spectrometers timing instruments, and motivates defining $\DP$ as
the primary signal primitive.

\subsection{Temporal programming}

Temporal programming is a paradigm, introduced for distributed
cyber-physical systems, in which the sole inter-node signal is a bare
pulse---a rising edge with no content, no header, no payload
\citep{kopetz2011,lamport1978}. The computational datum is
$\DP(k) = \Tref(k) - \trec(k)$: the signed difference between the
reference-oscillator-predicted arrival time of the $k$-th pulse and its
actual reception time. Programs are written as \emph{timing cells}:
measurable intervals in $\DP$-space. Execution dispatches a precompiled
action when the observed $\DP(k)$ falls in a cell.

No parser exists because there is no payload to parse. The security
consequence is absolute: buffer overflow, SQL injection, command
injection, and payload mutation are architecturally impossible---not
mitigated, but absent by construction \citep{saltzer1975}.

We show that this paradigm maps exactly onto MS acquisition, with the
result that a temporal programming MS acquisition system inherits the
full structural incorruptibility of the paradigm.

\subsection{Summary of contributions}

\begin{enumerate}[nosep,label=(\roman*)]
\item We derive all four mass analyzer equations from a single
  \emph{partition Lagrangian} (Section~\ref{sec:lagrangian}), establishing
  the unified physical basis.
\item We prove the \emph{Time-Count Identity}
  $dM/dt = \omega/(2\pi) = 1/\tp$ (Section~\ref{sec:timecount}), which
  identifies time as partition-state traversal.
\item We derive the \emph{Partition Uncertainty Relation}
  $\Delta M \cdot \tp \geq \hbar$ and use it to produce the
  \emph{real-time/resolution law} governing the minimum acquisition time
  for a given target cell (Section~\ref{sec:uncertainty}).
\item We define the temporal programming model for MS, prove the
  \emph{Structural Incorruptibility Theorem}, and give explicit compilers
  from $m/z$ targets to timing cells for all four analyzers
  (Sections~\ref{sec:tp}--\ref{sec:compiler}).
\item We prove the \emph{Composition-Inflation Theorem}
  $T(n,d) = d(d+1)^{n-1}$ and show that ion-mobility drift time provides
  a natural second timing axis, enabling isobar separation without
  additional mass resolution (Section~\ref{sec:multidim}).
\item We give a concrete implementation architecture for each instrument
  type (Section~\ref{sec:impl}).
\end{enumerate}

% =========================================================
\section{Physical Foundations}
\label{sec:physics}
% =========================================================

\subsection{Bounded phase space and oscillatory necessity}

\begin{axiom}[Bounded Phase Space]
\label{ax:bps}
Every persistent physical system occupies a bounded region
$\Omega \subset \R^{2d}$ of phase space with finite Liouville measure:
$\mu(\Omega) < \infty$.
\end{axiom}

\begin{theorem}[Liouville, \citealp{liouville1838}]
\label{thm:liouville}
Hamiltonian flow preserves the symplectic volume measure $\mu$:
for any measurable set $A \subset \Omega$ and any time $t$,
$\mu(\phi_t(A)) = \mu(A)$.
\end{theorem}

\begin{theorem}[Poincar{\'e} Recurrence, \citealp{poincare1890}]
\label{thm:poincare}
Let $\mu(\Omega) < \infty$ and let $\phi_t$ preserve $\mu$.
For any measurable $A \subset \Omega$ with $\mu(A) > 0$,
for $\mu$-almost every $x \in A$ there exists a sequence
$t_n \to \infty$ with $\phi_{t_n}(x) \in A$.
\end{theorem}

\begin{proof}
Set $A_k = \phi_{-kT}(A)$ for fixed $T > 0$. By Liouville,
$\mu(A_k) = \mu(A) > 0$. Since $\mu(\Omega) < \infty$ and the $A_k$
have constant positive measure, they cannot be pairwise disjoint.
There exist $j < k$ with $A_j \cap A_k \neq \emptyset$; any point in
this intersection returns to $A$ after $(k-j)T$ time units. Iterating
gives the required sequence. \qed
\end{proof}

\begin{corollary}[Oscillatory Necessity]
\label{cor:osc}
Every bounded Hamiltonian system without fixed points is oscillatory:
no coordinate can undergo monotonic escape, and every trajectory
returns to any neighbourhood of its initial condition.
\end{corollary}

\begin{proof}
Monotonic escape contradicts Theorem~\ref{thm:poincare} on a compact
surface. Absence of fixed points ($T > 0$ by the Third Law) forces
oscillation. \qed
\end{proof}

For integrable systems, the Arnold--Liouville theorem \citep{arnold1989}
guarantees action-angle coordinates $(J_i, \theta_i)$ with
$\dot{\theta}_i = \omega_i(J) = \partial H/\partial J_i = $ const.
Near-integrable systems are quasi-periodic on most tori by the KAM
theorem. The partition counting framework applies to the generic bounded
physical system.

\subsection{Partition coordinates and capacity}

\begin{definition}[Partition State]
\label{def:pstate}
A \emph{partition state} of a bounded three-dimensional Hamiltonian
system is an equivalence class of phase-space points sharing the same
closure label $(n, \ell, m, s)$ where:
$n \in \{1,2,3,\ldots\}$ (principal level, radial closure number),
$\ell \in \{0,1,\ldots,n-1\}$ (angular sublevel),
$m \in \{-\ell,\ldots,+\ell\}$ (azimuthal projection),
$s \in \{-\tfrac{1}{2}, +\tfrac{1}{2}\}$ (rotational parity).
The set of all partition states is $\calP$.
\end{definition}

\begin{theorem}[Partition Capacity]
\label{thm:capacity}
The number of distinct partition states at principal level $n$ is
$C(n) = 2n^2$.
\end{theorem}

\begin{proof}
For level $n$: $\ell$ ranges over $\{0,\ldots,n-1\}$; for each $\ell$,
$m$ ranges over $2\ell+1$ values; $s$ contributes a factor of 2:
\[
C(n) = 2\sum_{\ell=0}^{n-1}(2\ell+1) = 2\!\left(n + 2\cdot\tfrac{n(n-1)}{2}\right) = 2n^2.\qed
\]
\end{proof}

The cumulative count at or below level $n_{\max}$ is:
\begin{equation}
N_{\text{state}}(n_{\max}) = \sum_{n=1}^{n_{\max}} 2n^2
= \frac{n_{\max}(n_{\max}+1)(2n_{\max}+1)}{3}.
\label{eq:cumcount}
\end{equation}

\subsection{The partition Lagrangian}
\label{sec:lagrangian}

\begin{definition}[Partition Lagrangian]
\label{def:lagrangian}
The \emph{partition Lagrangian} governing a bounded oscillatory ion of
mass $m$ and charge $q$ in an electromagnetic environment is:
\begin{equation}
\LM = \tfrac{1}{2}\mu|\dot{\mathbf{x}}|^2
+ \mu\dot{\mathbf{x}}\cdot\AM(\mathbf{x},t)
- \Mop(\mathbf{x},t),
\label{eq:lagrangian}
\end{equation}
where $\mu = \alpha(m/z)$ is the partition inertia ($z = q/e$ the charge
state, $\alpha$ an instrument calibration constant fixed by one reference
measurement), $\AM$ is the partition vector potential encoding the
confining electromagnetic field, and $\Mop$ is the partition Hamiltonian
density.
\end{definition}

\begin{theorem}[Lorentz Force as Euler-Lagrange Equation]
\label{thm:lorentz}
With $\Mop = q\phi$ (electrostatic potential) and
$\AM = (q/\mu)\mathbf{A}$ (magnetic vector potential),
the Euler-Lagrange equations of $\LM$ give the Lorentz force law:
\begin{equation}
\mu\ddot{\mathbf{x}} = q\mathbf{E} + q\dot{\mathbf{x}}\times\mathbf{B},
\label{eq:lorentz}
\end{equation}
where $\mathbf{E} = -\nabla\phi - \partial\mathbf{A}/\partial t$ and
$\mathbf{B} = \nabla\times\mathbf{A}$.
\end{theorem}

\begin{proof}
The Euler-Lagrange equation in component $i$:
\[
\frac{d}{dt}\frac{\partial\LM}{\partial\dot{x}_i} - \frac{\partial\LM}{\partial x_i} = 0.
\]
Computing each term with $A_i^{\mathcal{M}} = (q/\mu)A_i$:
\[
\frac{\partial\LM}{\partial\dot{x}_i} = \mu\dot{x}_i + \mu A_i^{\mathcal{M}},
\]
\[
\frac{d}{dt}(\mu\dot{x}_i + \mu A_i^{\mathcal{M}}) =
\mu\ddot{x}_i + \mu(\partial_t A_i^{\mathcal{M}} + \dot{x}_j\partial_j A_i^{\mathcal{M}}),
\]
\[
\frac{\partial\LM}{\partial x_i} = \mu\dot{x}_j\partial_i A_j^{\mathcal{M}} - \partial_i\Mop.
\]
Subtracting and rearranging:
\[
\mu\ddot{x}_i = -\mu\partial_t A_i^{\mathcal{M}} + \mu\dot{x}_j(\partial_i A_j^{\mathcal{M}} - \partial_j A_i^{\mathcal{M}}) - \partial_i\Mop.
\]
Substituting $A_i^{\mathcal{M}} = (q/\mu)A_i$ and $\Mop = q\phi$:
\[
\mu\ddot{x}_i = -q\partial_t A_i + q\dot{x}_j(\partial_i A_j - \partial_j A_i) - q\partial_i\phi.
\]
The first and last terms give $qE_i$; the middle term gives $q(\dot{\mathbf{x}}\times\mathbf{B})_i$. \qed
\end{proof}

% =========================================================
\section{All Four Analyzer Equations from One Lagrangian}
\label{sec:analyzers}
% =========================================================

The four commercially deployed mass analyzer equations are obtained
from~\eqref{eq:lagrangian} by choosing specific forms of $\AM$ and $\Mop$.
No physical assumption beyond the Lagrangian and the calibration of
the field is required.

\subsection{Time-of-flight analyzer}

\begin{theorem}[TOF Equation]
\label{thm:tof}
For uniform acceleration $\Mop = qV$ (constant accelerating potential),
$\AM = 0$, and field-free drift of length $L$, the Euler-Lagrange
equations of $\LM$ give:
\begin{equation}
T_{\mathrm{TOF}} = L\sqrt{\frac{m}{2qV}}.
\label{eq:tof}
\end{equation}
\end{theorem}

\begin{proof}
In the acceleration gap, $\mu\ddot{x} = -\partial\Mop/\partial x = 0$
(uniform field: $qV$ is constant in $x$). Energy conservation from
rest: $\tfrac{1}{2}\mu v^2 = qV$, giving $v = \sqrt{2qV/\mu}$.
With $\mu = \alpha m/z$ and $\alpha = z$ (charge calibration):
$v = \sqrt{2qV/m}$.
In the field-free drift: $T = L/v = L\sqrt{m/(2qV)}$. \qed
\end{proof}

\subsection{Quadrupole analyzer}

\begin{theorem}[Mathieu Stability]
\label{thm:mathieu}
For the quadrupole potential
$\Mop = \tfrac{q}{r_0^2}(U + V_0\cos\Omega t)(x^2 - y^2)$
and $\AM = 0$, the Euler-Lagrange equations reduce to the
Mathieu equations:
\begin{equation}
\frac{d^2u}{d\xi^2} + (a_u - 2q_u\cos 2\xi)u = 0,\quad
\xi = \tfrac{1}{2}\Omega t,
\label{eq:mathieu}
\end{equation}
with Mathieu parameters $a_u = \pm 4qU/(m\Omega^2 r_0^2)$ and
$q_u = \pm 2qV_0/(m\Omega^2 r_0^2)$. Ion transmission requires
bounded solutions, which occur in the first stability region:
$0 < a_u < 0.237$, $0 < q_u < 0.908$; the mass scan line is
$a_z = -2q_z$ \citep{march1989}.
\end{theorem}

\begin{proof}
The $x$-component Euler-Lagrange equation:
$\mu\ddot{x} = -\partial\Mop/\partial x = -2q(U+V_0\cos\Omega t)x/r_0^2$.
With $\mu = m$ and $\xi = \Omega t/2$:
$d^2x/d\xi^2 + (a_x - 2q_x\cos 2\xi)x = 0$,
where $a_x = -4qU/(m\Omega^2 r_0^2)$ and $q_x = 2qV_0/(m\Omega^2 r_0^2)$.
The $y$-equation has opposite signs for $a$ and $q$. Stability in both
$x$ and $y$ simultaneously defines the first stability region \citep{mathieu1868,march1989}.
\qed
\end{proof}

\subsection{Orbitrap analyzer}

\begin{theorem}[Orbitrap Frequency]
\label{thm:orbitrap}
For the harmonic axial potential $\Mop = \tfrac{1}{2}q\kappa z^2$
and an azimuthal vector potential $\AM = A_\phi\hat{\phi}$,
the Euler-Lagrange equation in $z$ gives axial harmonic oscillation at:
\begin{equation}
\omega_z = \sqrt{\frac{q\kappa}{m}}.
\label{eq:orbitrap}
\end{equation}
The $m/z$ ratio is recovered as $m/z = e\kappa/\omega_z^2$
\citep{makarov2000}.
\end{theorem}

\begin{proof}
The $z$-component: $\mu\ddot{z} = -\partial\Mop/\partial z = -q\kappa z$.
With $\mu = m$: $m\ddot{z} = -q\kappa z$, giving
$\omega_z^2 = q\kappa/m$. \qed
\end{proof}

\subsection{FT-ICR analyzer}

\begin{theorem}[Cyclotron Frequency]
\label{thm:fticr}
For a uniform magnetic field $\mathbf{B} = B\hat{z}$ encoded as
$\AM = (q/\mu)(B/2)(-y, x, 0)$ and $\Mop = 0$, the Euler-Lagrange
equations give circular motion at the cyclotron frequency:
\begin{equation}
\omega_c = \frac{qB}{m} = \frac{eB}{m/z}.
\label{eq:fticr}
\end{equation}
The $m/z$ ratio is recovered as $m/z = eB/\omega_c$
\citep{marshall1998}.
\end{theorem}

\begin{proof}
With $A_x^{\mathcal{M}} = -(q/\mu)(B/2)y$ and
$A_y^{\mathcal{M}} = (q/\mu)(B/2)x$:
\[
\mu\ddot{x} = \mu\dot{y}\!\left(\frac{\partial A_y^{\mathcal{M}}}{\partial x}
- \frac{\partial A_x^{\mathcal{M}}}{\partial y}\right) = q\dot{y}B,
\quad
\mu\ddot{y} = -q\dot{x}B.
\]
Writing $\zeta = x + iy$: $\ddot{\zeta} = -i(qB/\mu)\dot{\zeta}$,
giving $\dot{\zeta}(t) = \dot{\zeta}(0)e^{-i\omega_c t}$ with
$\omega_c = qB/\mu = zeBz/m = eB z/m$. \qed
\end{proof}

\begin{remark}[Four analyzers, one Lagrangian]
Theorems~\ref{thm:tof}--\ref{thm:fticr} demonstrate that TOF, quadrupole,
Orbitrap, and FT-ICR are four boundary conditions on $\LM$. Their
historically separate derivations and design frameworks are unified
under the partition Lagrangian. All four measure $m/z$ via a timing
or frequency observable; none measures $m/z$ directly.
\end{remark}

% =========================================================
\section{Time as Partition-State Traversal}
\label{sec:timecount}
% =========================================================

\subsection{The Time-Count Identity}

The identification of time with partition-state counting is not an
approximation; it is exact for any periodic oscillator.

\begin{definition}[Reference Oscillator]
\label{def:oscillator}
A \emph{reference oscillator} is a tuple $(\fref, \epsilon, M_0)$ where:
$\fref > 0$ is the nominal frequency,
$\epsilon \in [0,1)$ is the maximum fractional frequency deviation,
and $M_0 \in \N$ is the initial cycle count.
The oscillator generates reference times $\Tref(k) = k/\fref$ for
$k \geq M_0$, and the cycle count at physical time $t$ is
$\Mcyc(t) = M_0 + \lfloor(t - \Tref(M_0))\fref\rfloor$.
\end{definition}

\begin{theorem}[Time-Count Identity]
\label{thm:timecount}
For a bounded oscillatory system with angular frequency $\omega$ and
mean partition residence time $\tp = 2\pi/\omega$:
\begin{equation}
\frac{d\Mcyc}{dt} = \frac{\omega}{2\pi} = \frac{1}{\tp}.
\label{eq:timecount}
\end{equation}
This identity is exact: $\Mcyc(t) = f_{\mathrm{osc}} \cdot t + \Mcyc(0)$
where $f_{\mathrm{osc}} = \omega/(2\pi)$.
\end{theorem}

\begin{proof}
The oscillator completes one cycle per period $T_p = 2\pi/\omega$, so
$d\Mcyc/dt = 1/T_p = \omega/(2\pi)$. Integrating with the definition
$\tp = T_p$: $\Mcyc(t) = t/\tp + \Mcyc(0)$. The equality
$\omega/(2\pi) = 1/\tp$ follows directly from $T_p = 2\pi/\omega$. \qed
\end{proof}

\begin{corollary}[Strict Monotonicity]
\label{cor:mono}
For $t_2 > t_1$, $\Mcyc(t_2) > \Mcyc(t_1)$: the partition count is
strictly non-decreasing. No physical operation can decrement $M$.
\end{corollary}

\begin{proof}
$\Mcyc(t_2) - \Mcyc(t_1) = f_{\mathrm{osc}}(t_2 - t_1) > 0$ for
$t_2 > t_1$ and $f_{\mathrm{osc}} > 0$. \qed
\end{proof}

The corollary has a direct security consequence for MS acquisition:
a recorded timing deviation from cycle $M_0$, if replayed at cycle
$M_1 > M_0$, produces a shifted deviation
$\DP_1 = \DP_0 + (M_1-M_0)/\fref \neq \DP_0$, placing it in a different
cell or outside the calibrated window. Replay attacks are intrinsically
neutralised by the monotone count.

\subsection{MS observables in terms of the partition count}

The four analyzer frequencies in terms of the cycle count:

\begin{align}
\omega_{\mathrm{TOF}} &= \frac{2\pi}{T_{\mathrm{TOF}}} = \frac{2\pi}{L}\sqrt{\frac{2qV}{m}},
\label{eq:tof-omega}\\
\omega_z^{\mathrm{Orb}} &= \sqrt{\frac{q\kappa}{m}},
\label{eq:orb-omega}\\
\omega_c^{\mathrm{ICR}} &= \frac{qB}{m},
\label{eq:icr-omega}\\
\omega_{\mathrm{quad}} &= \Omega\sqrt{\tfrac{1}{2}q_u} \text{ (secular frequency)}.
\label{eq:quad-omega}
\end{align}

In each case, $dM/dt = \omega/(2\pi)$. The $m/z$ value governs the rate
at which the instrument's partition counter advances. A lighter ion
advances the counter faster (higher $\omega$); a heavier ion advances it
slower.

% =========================================================
\section{The Partition Uncertainty and the Resolution Law}
\label{sec:uncertainty}
% =========================================================

\subsection{Derivation of the partition uncertainty relation}

\begin{theorem}[Partition Uncertainty Relation]
\label{thm:uncertainty}
For a partition state with Hamiltonian density $\Mop$ (playing the role
of energy) and mean residence time $\tp$ (playing the role of time):
\begin{equation}
\Delta\Mop \cdot \tp \geq \hbar,
\label{eq:uncertainty}
\end{equation}
where $\Delta\Mop$ is the standard deviation of $\Mop$ over the
partition ensemble.
\end{theorem}

\begin{proof}
This is the Robertson--Schr{\"o}dinger time-energy uncertainty relation
\citep{robertson1929,heisenberg1927} applied to the partition basis.
The observable $\hat{\Mop}$ is conjugate to the partition time $\hat{t}_p$
via $[\hat{\Mop}, \hat{t}_p] = i\hbar$ (canonical conjugacy of energy
and time). The Robertson bound then gives
$\Delta\Mop \cdot \Delta t_p \geq \tfrac{1}{2}|{\langle[\hat{\Mop},\hat{t}_p]\rangle}|
= \hbar/2$.
With $\Delta t_p = \tp$ as the natural timescale:
$\Delta\Mop \cdot \tp \geq \hbar$. \qed
\end{proof}

\subsection{The real-time/resolution law}

\begin{theorem}[Real-Time Resolution Law]
\label{thm:resolution-law}
For a target cell $\calC_i$ of width $\delta_{m/z}$ in mass-to-charge
space, the corresponding width in partition-count space is:
\begin{equation}
\delta M = \frac{\omega}{\pi}\cdot\delta t,
\label{eq:dm-dt}
\end{equation}
where $\delta t$ is the timing resolution.
The minimum acquisition time $\tau_{\min}$ before the trigger can
reliably assign cell membership is:
\begin{equation}
\tau_{\min}(\calC_i) = \frac{\hbar}{\delta\Mop_i},
\label{eq:taumin}
\end{equation}
where $\delta\Mop_i$ is the partition energy width of cell $\calC_i$.
\end{theorem}

\begin{proof}
From the uncertainty relation~\eqref{eq:uncertainty}:
$\Delta\Mop \cdot \tau \geq \hbar$, so
$\tau \geq \hbar/\Delta\Mop$. The minimum acquisition time is achieved
when $\Delta\Mop = \delta\Mop_i$ (the cell width in partition energy
units), giving $\tau_{\min} = \hbar/\delta\Mop_i$. A narrower cell
(smaller $\delta\Mop_i$) requires a longer minimum observation time.
\qed
\end{proof}

\begin{corollary}[Dynamic Acquisition Schedule]
\label{cor:dynamic}
For a target list $\{\calC_1, \ldots, \calC_m\}$ with cell widths
$\delta\Mop_1 \leq \ldots \leq \delta\Mop_m$, the acquisition should
allocate time $\tau_j \geq \hbar/\delta\Mop_j$ to each target $j$.
Wide cells (low-specificity targets) allow earlier triggering; narrow
cells (high-specificity targets) require longer transients. This
dynamic schedule is set by equation~\eqref{eq:taumin} from first
principles, without empirical calibration.
\end{corollary}

Table~\ref{tab:resolution} shows the partition-uncertainty-derived minimum
acquisition times for typical Orbitrap target windows at 10\,MHz
reference frequency.

\begin{table}[h]
\centering
\caption{Minimum acquisition times from the partition uncertainty relation
for Orbitrap acquisition at $f_{\text{ref}} = 10$\,MHz.}
\label{tab:resolution}
\small
\begin{tabular}{lcc}
\toprule
Target window $\delta_{m/z}$ & $\delta\Mop_i$ (arb.) & $\tau_{\min}$ \\
\midrule
10\,Da (survey) & wide & $\sim$1\,ms \\
1\,Da (precursor) & medium & $\sim$10\,ms \\
0.1\,Da (PRM) & narrow & $\sim$100\,ms \\
0.01\,Da (targeted) & very narrow & $\sim$1\,s \\
\bottomrule
\end{tabular}
\end{table}

The partition uncertainty relation thus provides the exact design curve
for the specificity-speed tradeoff that is currently set empirically in
all commercial acquisition software.

% =========================================================
\section{The Temporal Programming Model}
\label{sec:tp}
% =========================================================

\subsection{The timing deviation primitive}

\begin{definition}[Timing Deviation]
\label{def:dp}
Let $(\fref, \epsilon, M_0)$ be a reference oscillator. For the $k$-th
detected ion or signal with actual reception time $\trec(k)$, the
\emph{timing deviation} is:
\begin{equation}
\DP(k) = \Tref(k) - \trec(k) \in \R.
\label{eq:dp}
\end{equation}
$\DP(k) > 0$ indicates early arrival; $\DP(k) < 0$ indicates late arrival;
$\DP(k) = 0$ indicates perfect agreement with the reference.
\end{definition}

In a TOF instrument, $\Tref(k)$ is the target flight time of the $k$-th
target ion computed from the cell compiler (Section~\ref{sec:compiler}),
and $\trec(k)$ is the measured flight time. In an Orbitrap, $\Tref(k)$
is the target oscillation phase at transient sample $k$, and $\trec(k)$
is the measured phase. In FT-ICR, $\Tref(k)$ is the reference cyclotron
phase.

\subsection{Timing cells}

\begin{definition}[Timing Cell]
\label{def:cell}
A \emph{timing cell} is a Borel-measurable set $\calC \subseteq \calD$
with positive Lebesgue measure $\Leb(\calC) > 0$, where
$\calD \subset \R$ is the calibrated $\DP$-space of the instrument.
A \emph{cell partition} of $\calD$ is a finite collection
$\{\calC_1, \ldots, \calC_m\}$ satisfying:
(i) $\bigcup_i \calC_i = \calD$ (coverage),
(ii) $\calC_i \cap \calC_j = \emptyset$ for $i \neq j$ (disjointness),
(iii) $\Leb(\calC_i) > 0$ for all $i$ (non-degeneracy).
\end{definition}

The positive-measure requirement is not cosmetic. A point predicate
$\{\DP = c\}$ has Lebesgue measure zero and cannot be robustly matched
in the presence of oscillator jitter. Physical cells must have positive
width, bounded below by the instrument's timing uncertainty.

\subsection{Cell registry and action dispatch}

\begin{definition}[Cell Registry and Action Map]
A \emph{cell registry} $\Gamma$ is a pair $(\{\calC_i\}, \calA)$ where
$\{\calC_i\}$ is a cell partition of $\calD$ and
$\calA: \{\calC_i\} \to \mathcal{F}$ maps each cell to a precompiled
action (e.g.\ ACCEPT, DISCARD, FRAGMENT, RECORD). The registry is
\emph{frozen at compile time} and stored in read-only memory.
\end{definition}

\subsection{Operational semantics}

A \emph{configuration} is a tuple $(\sigma, \Gamma, \tau, \phi)$ where:
$\sigma \in \{\mathrm{COMPILE},\, \mathrm{EXECUTE},\, \mathrm{IDLE}\}$,
$\Gamma$ is the frozen cell registry,
$\tau$ is the current trajectory (sequence of accumulated $\DP$ values),
and $\phi$ is the pending action queue.

The reduction rules are:

\vspace{4pt}
\noindent\textbf{(Receive):} While $|\tau| < n$ and $\sigma = \mathrm{COMPILE}$,
append each arriving $\DP(k)$ to $\tau$.

\noindent\textbf{(Assign):} When $|\tau| = n$, look up $\tau$ in $\Gamma$;
if found, transition to $\mathrm{EXECUTE}$ with action $\calA(\calC)$.

\noindent\textbf{(Fire):} In $\mathrm{EXECUTE}$, dispatch each action in $\phi$.

\noindent\textbf{(Reset):} When $\phi$ is empty, return to $\mathrm{IDLE}$;
then immediately to $\mathrm{COMPILE}$.
\vspace{4pt}

\begin{theorem}[Phase Mutual Exclusion]
\label{thm:phase}
In any execution trace, COMPILE and EXECUTE never hold simultaneously.
\end{theorem}

\begin{proof}
Each configuration carries exactly one phase tag. No reduction rule
produces two simultaneous tags. COMPILE requires an open trajectory
($|\tau| < n$, cell unknown); EXECUTE requires a closed trajectory
($|\tau| = n$, cell known). These preconditions are contradictory. \qed
\end{proof}

% =========================================================
\section{Structural Incorruptibility}
\label{sec:incorruptibility}
% =========================================================

\subsection{The adversary model}

Consider an adversary $\mathcal{ADV}$ with capabilities:
\textbf{C1} (content control): choose arbitrary bit patterns in any signal;
\textbf{C2} (timing control): shift arrival times by $\pm\Delta_{\max}$;
\textbf{C3} (injection): insert additional pulses;
\textbf{C4} (deletion): suppress pulses;
\textbf{C5} (replay): retransmit recorded signals.

\subsection{Main theorem}

\begin{theorem}[Structural Incorruptibility]
\label{thm:incorruptibility}
In a temporal programming MS acquisition system, capability \textbf{C1}
(content control) provides zero leverage: no content manipulation can
cause the trigger to accept an ion that the timing deviations would
not accept.
\end{theorem}

\begin{proof}
A received ion event has two measurable properties: the channel
(the physical flight path or detector electrode) and $\DP(k)$. The
channel is determined by the physical geometry of the instrument. The
$\DP(k)$ value depends only on the arrival time, not on the ion's
content (charge distribution, internal energy, or isotopic pattern).

Signal processing consists only of the comparison
$\DP(k) \in \calC_i ?$ for each cell in the registry. This comparison
involves no parsing, no buffer operations, and no code interpretation.
The cell-action map $\calA$ is stored in ROM and no runtime path leads
from signal content to $\calA$ modification.

Therefore, the set of executable actions is always
$\{\calA(\calC_i): \calC_i \in \Gamma\}$, fixed at compile time.
$\mathcal{ADV}$ cannot expand this set using \textbf{C1}. \qed
\end{proof}

\begin{corollary}[Off-Target Ions Are Structurally Rejected]
\label{cor:offtarget}
An off-target ion --- one whose $m/z$ was not compiled into $\Gamma$ ---
arrives with $\DP(k)$ drawn from a distribution centred away from every
cell in $\Gamma$. By the calibration of the cell compiler
(Section~\ref{sec:compiler}), its $\DP(k)$ falls outside all cells.
Rejection is structural, not statistical: the ion does not fail a
threshold test; it simply does not fall in any cell.
\end{corollary}

This is the mass spectrometry analogue of the statement that selection
rules are \emph{dead zones}, not rare events. A symmetry-forbidden
transition does not occur at low probability; it does not occur.
Similarly, an off-target ion is not accepted at low rate; it is not
accepted.

\begin{corollary}[Replay Attacks Are Neutralised]
\label{cor:replay}
An adversary recording $\DP_0 = \DP(k)$ at cycle $M_0$ and replaying
it at cycle $M_1 > M_0$ produces:
$\DP_1 = \DP_0 + (M_1 - M_0)/\fref$.
Since $M_1 > M_0$ and $\fref > 0$, the shift $\delta = (M_1-M_0)/\fref > 0$
displaces the replayed signal from its original cell by the amount
$\delta$. For $\delta \geq \min_i \Leb(\calC_i)$ (cell width), the
replay lands in a different cell or outside $\calD$. As $M$ grows
monotonically (Corollary~\ref{cor:mono}), $\delta$ grows without
bound, and replay eventually fails structurally.
\end{corollary}

% =========================================================
\section{Compiling Targets to Timing Cells}
\label{sec:compiler}
% =========================================================

The compiler translates an $m/z$ target window $[m_0/z, m_1/z]$ into
a timing cell $[a_i, b_i] \subset \calD$ for each analyzer type. This
compilation is exact: once compiled, the registry requires no $m/z$
arithmetic during acquisition.

\subsection{TOF compiler}

From Theorem~\ref{thm:tof}, a target in the window
$[m_0/z, m_1/z]$ maps to flight times:
\begin{equation}
[T_0, T_1] = \left[L\sqrt{\frac{m_0}{2qV}},\; L\sqrt{\frac{m_1}{2qV}}\right].
\label{eq:tof-cell}
\end{equation}
The timing cell is $\calC_i = [T_{\mathrm{ref}} - T_1, T_{\mathrm{ref}} - T_0]$
in $\DP$-space, where $T_{\mathrm{ref}}$ is the reference flight time
for the calibrant ion.

\subsection{Orbitrap compiler}

From Theorem~\ref{thm:orbitrap}, the target $[m_0/z, m_1/z]$ maps to
frequencies:
\begin{equation}
[\omega_1, \omega_0] = \left[\sqrt{\frac{q\kappa}{m_1/z}},\;\sqrt{\frac{q\kappa}{m_0/z}}\right],
\label{eq:orb-cell}
\end{equation}
(note inversion: higher mass → lower frequency). In $\DP$-space, this
corresponds to a phase-difference window computed from partial-transient
frequency estimation.

\subsection{FT-ICR compiler}

From Theorem~\ref{thm:fticr}:
\begin{equation}
[\omega_{c,1}, \omega_{c,0}] = \left[\frac{eB}{m_1/z},\; \frac{eB}{m_0/z}\right].
\label{eq:icr-cell}
\end{equation}
Again the mapping inverts: lighter ions cycle faster. The timing cell
in $\DP$-space is a phase interval computed from the reference cyclotron
period.

\subsection{Quadrupole compiler}

From Theorem~\ref{thm:mathieu}, the scan line $a_z = -2q_z$ at mass
$m$ selects $q_z = 2qV_0/(m\Omega^2 r_0^2)$. For target window
$[m_0/z, m_1/z]$:
\begin{equation}
[q_{z,\min}, q_{z,\max}] = \left[\frac{2qV_0}{(m_1/z)\Omega^2 r_0^2},\;
\frac{2qV_0}{(m_0/z)\Omega^2 r_0^2}\right].
\label{eq:quad-cell}
\end{equation}
For a fixed RF voltage $V_0$ and frequency $\Omega$, this is a window
on the dimensionless Mathieu parameter, which maps linearly to the
transmission timing of the stability zone.

\begin{table}[ht]
\centering
\caption{The temporal programming map: mass spectrometry to timing
programming primitives.}
\label{tab:mapping}
\small
\begin{tabular}{p{3.8cm}p{3.8cm}}
\toprule
\textbf{Temporal programming} & \textbf{Mass spectrometry} \\
\midrule
Reference oscillator $\fref$ & TOF start pulse / Orbitrap oscillator / ICR reference \\
$\trec(k)$ & Ion arrival time (TOF) or transient phase (Orb/ICR) \\
$\DP(k)$ & Timing/frequency deviation from calibrated target \\
Timing cell $\calC_i$ & Target $m/z$ window in time space \\
Cell registry $\Gamma$ & SRM/PRM inclusion list, compiled to timing \\
Dispatch $\calA$ & Accept (record/fragment) or discard, real-time \\
COMPILE phase & Transient/flight accumulation \\
EXECUTE phase & Trigger decision and action dispatch \\
Replay immunity & Monotone cycle count rejects recorded signals \\
Structural incorruptibility & No $m/z$ parsing: no off-target admission \\
\bottomrule
\end{tabular}
\end{table}

% =========================================================
\section{Multi-Dimensional Acquisition and Isobar Separation}
\label{sec:multidim}
% =========================================================

\subsection{The isobar limitation}

A single timing axis resolves $m/z$ but cannot separate two species at
identical $m/z$ (isobars). Both arrive with the same $\DP$ and land in
the same cell. This is a genuine limitation of one-dimensional temporal
programming, not a detail to be finessed.

The resolution: add a second timing axis. Ion-mobility spectrometry
(IMS) separates ions by their collision cross-section (CCS) in a drift
gas \citep{clemmer1997,lanucara2014}. The drift time $t_d \propto$
CCS/$\sqrt{m/z}$ is, again, a timing observable. Isobars that share
$m/z$ but differ in CCS produce different drift times and different
$\DP_d$ values on the second axis.

\subsection{Composition-inflation theorem}

\begin{theorem}[Composition-Inflation]
\label{thm:inflation}
For a temporal programming system with $d$ independent timing channels
and $n$ timing events, the number of categorically distinguishable
labeled trajectories is:
\begin{equation}
T(n, d) = d(d+1)^{n-1}.
\label{eq:inflation}
\end{equation}
\end{theorem}

\begin{proof}
A labeled composition of $n$ into $k$ parts assigns each part a label
in $\{1,\ldots,d\}$. The number of compositions of $n$ into $k$ parts
is $\binom{n-1}{k-1}$ (stars and bars \citep{andrews1976}).
For each $k$-part composition, there are $d^k$ label assignments.
Summing over $k$:
\[
T(n,d) = \sum_{k=1}^{n}\binom{n-1}{k-1}d^k
= d\sum_{j=0}^{n-1}\binom{n-1}{j}d^j = d(1+d)^{n-1},
\]
where the last step is the binomial theorem. \qed
\end{proof}

\begin{corollary}[IMS-MS Capacity]
For a two-channel IMS-MS system ($d = 2$):
$T(n, 2) = 2 \cdot 3^{n-1}$.
At $n = 10$ timing events: $T(10, 2) = 2 \cdot 3^9 = 39{,}366$
distinguishable ion categories---orders of magnitude more than any
practical target list. Two timing channels resolve all isobars
distinguishable by CCS.
\end{corollary}

\begin{corollary}[Growth Properties]
(i) $T(1,d) = d$: one event, $d$ categories.
(ii) $T(n+1, d) = (d+1) \cdot T(n,d)$: each additional event multiplies
the state count by $(d+1)$.
(iii) $T(n,d) \gg n \cdot d$ for large $n$: the multi-channel state space
grows exponentially faster than any linear estimate.
\end{corollary}

\subsection{Two-dimensional cell registry}

For an IMS-MS system, the cell registry is a 2D partition of
$(\DP_{\mathrm{flight}}, \DP_{\mathrm{drift}}) \in \calD_1 \times \calD_2$.
Each 2D cell $\calC_{ij} = \calC_i^{(1)} \times \calC_j^{(2)}$ admits
ions within a specific $m/z$ window AND a specific CCS window.
The compilation maps $(m/z, \mathrm{CCS})$ targets to product cells.
This separates isobars with different CCS by assigning them to different
cells: they are structurally separated, not statistically filtered.

\begin{theorem}[Isobar Separation by Two-Channel Timing]
Two ion species $A$ and $B$ with $m/z_A = m/z_B$ but
$\mathrm{CCS}_A \neq \mathrm{CCS}_B$ produce:
$\DP_{\mathrm{drift},A} \neq \DP_{\mathrm{drift},B}$.
Therefore, for cells $\calC^{(2)}_A$ and $\calC^{(2)}_B$ compiled
from their respective CCS values:
$A \in \calC_{i,A}$ and $B \in \calC_{i,B}$ with $\calC_{i,A} \neq \calC_{i,B}$.
The two species land in structurally distinct cells.
\end{theorem}

\begin{proof}
The drift time $t_{d,X} = \ell/(K_0 E) \cdot \sqrt{m_X/z_X}$, where
$K_0$ is the reduced mobility and $E$ the electric field. For
$m/z_A = m/z_B$ but $K_{0,A} \neq K_{0,B}$ (different CCS implies
different mobility): $t_{d,A} \neq t_{d,B}$, hence
$\DP_{d,A} \neq \DP_{d,B}$, and with calibrated cell widths less than
$|t_{d,A} - t_{d,B}|$, they fall in distinct cells. \qed
\end{proof}

% =========================================================
\section{Implementation}
\label{sec:impl}
% =========================================================

The implementation specifics differ by instrument type because the
actuator that responds to the trigger decision operates at different
points in the ion's trajectory.

\subsection{Q-TOF (quadrupole time-of-flight)}

In a Q-TOF, the quadrupole precedes the TOF region. The quadrupole
operates as a mass filter: for a target in $\Gamma$, the quadrupole
parameters ($U$, $V_0$) are compiled from the cell to admit only that
$m/z$ range. The temporal programming trigger drives the quadrupole
voltage schedule: each cell in $\Gamma$ corresponds to a compiled
$(a_u, q_u)$ pair. The TOF's reference oscillator provides $\Tref(k)$;
the flight time $\trec(k)$ is latched at ion arrival by an interrupt
service routine (ISR), and $\DP(k)$ is computed in $\mathcal{O}(1)$.

\begin{algorithm}[h]
\caption{Q-TOF temporal programming acquisition}
\label{alg:qtof}
\begin{algorithmic}[1]
\State Load compiled registry $\Gamma$ from ROM
\For{each target window $\calC_i \in \Gamma$}
  \State Set quadrupole $(U_i, V_{0,i})$ from compiler map~\eqref{eq:quad-cell}
  \State Trigger TOF pulse
  \State Latch $\trec$ at ion arrival via ISR
  \State Compute $\DP \gets \Tref - \trec$
  \If{$\DP \in \calC_i^{\mathrm{TOF}}$} ACCEPT \Else\ DISCARD
  \EndIf
\EndFor
\end{algorithmic}
\end{algorithm}

\subsection{Orbitrap}

The Orbitrap trigger operates on the partial transient. For target
cell $\calC_i$ with width $\delta\Mop_i$, the minimum transient
length is $\tau_{\min} = \hbar/\delta\Mop_i$ (Theorem~\ref{thm:resolution-law}).
A sliding-window frequency estimator (e.g.\ STFT or filter bank) runs
on the incoming image current. Once $\tau_{\min}$ samples are
accumulated, the estimated frequency $\hat\omega$ is compared to the
compiled target: $\DP \gets |\hat\omega - \omega_{\mathrm{ref}}|$.
If $\DP < \delta\omega_i/2$, the event is accepted.

\begin{algorithm}[h]
\caption{Orbitrap partial-transient trigger}
\label{alg:orbitrap}
\begin{algorithmic}[1]
\State Load $\Gamma = \{(\omega_i, \delta\omega_i)\}$ from ROM
\State $\tau \gets 0$; accumulate image current samples
\For{each sample}
  \State $\tau \gets \tau + 1/f_{\mathrm{ADC}}$
  \If{$\tau \geq \tau_{\min}(\calC_i)$}
    \State $\hat\omega \gets$ STFT frequency estimate
    \State $\DP \gets |\hat\omega - \omega_i|$
    \If{$\DP \leq \delta\omega_i/2$} TRIGGER ACCEPT; break
    \EndIf
  \EndIf
\EndFor
\end{algorithmic}
\end{algorithm}

The key property: the trigger is issued with the minimum necessary
transient, not a fixed preset transient. Narrow-window (high-resolution)
targets accumulate longer; wide-window (survey) targets trigger faster.

\subsection{FT-ICR}

The FT-ICR trigger is identical in structure to the Orbitrap trigger,
with $\omega_c = eB/(m/z)$ replacing $\omega_z = \sqrt{e\kappa/(m/z)}$.
The drift time of the cell registry scales as $m/z$ (linear) rather
than $\sqrt{m/z}$ (Orbitrap), which means FT-ICR cells have a
different shape in $\DP$-space. The minimum transient is the same:
$\tau_{\min} = \hbar/\delta\Mop_i$.

\subsection{Minimum cell width from oscillator stability}

The minimum physical cell width is bounded below by oscillator stability
and propagation jitter:
\begin{equation}
w_{\min} \geq \frac{2\epsilon}{\fref} + 2\delta_{\max},
\label{eq:cellwidth}
\end{equation}
where $\epsilon$ is the fractional frequency stability and
$\delta_{\max}$ is the maximum timing jitter. For a 10\,MHz OCXO
($\epsilon \approx 10^{-8}$): $w_{\min} \approx 2\,\mathrm{ns}$,
corresponding to $m/z$ resolution better than 1\,Da at 100\,Da.

% =========================================================
\section{Experimental Predictions and Validation}
\label{sec:validation}
% =========================================================

The temporal programming framework makes quantitative predictions
testable against existing and future MS data.

\subsection{The resolution-time tradeoff curve}

Prediction~1: For Orbitrap measurements of the same compound with
varying transient lengths $\tau$, the observed mass accuracy should
follow $\sigma(m/z) \propto 1/\tau$ (from
Theorem~\ref{thm:resolution-law}). This is already known empirically
\citep{hu2005} but has not been derived from first principles. The
partition uncertainty derivation provides both the functional form and
the proportionality constant in terms of the physical parameters
$\kappa$ and $q$.

\subsection{Off-target rejection rates}

Prediction~2: For a temporal programming acquisition with cell width
$w_{\calC}$, the fraction of chemical noise ions that land in any cell
is $\leq w_{\calC}/\calD$ (the ratio of cell width to $\DP$-space
width). For a 10-target acquisition over a 1000-Da $m/z$ range with
individual windows of 0.1\,Da, this is $10 \times 0.1/1000 = 10^{-3}$.
The rejection rate is structural and quantitatively predictable,
independent of signal intensity.

\subsection{IMS-MS isobar separation}

Prediction~3: For isobars with known CCS values
$\Omega_A \neq \Omega_B$, the drift time difference
$\Delta t_d = |\Omega_A - \Omega_B| \cdot \ell/(K_0 E)$ is computable
from the mobility equation. For $|\Omega_A - \Omega_B| > 5\,\text{\AA}^2$
at typical TIMS or DTIMS conditions, the species fall in distinct cells.
Prediction: the cell-based isobar separation rate should equal the
fraction of isobaric pairs with $|\Omega_A - \Omega_B| > w_{\calC_d}$
(drift cell width in CCS units).

\subsection{Replay immunity experimental test}

Prediction~4: A recorded ion timing signature from cycle $M_0$, when
artificially replayed at cycle $M_1 > M_0$ (achievable in a controlled
lab setting by injecting a synthetic signal), should produce
$\DP_1 = \DP_0 + (M_1-M_0)/\fref$. For a 10\,MHz oscillator and
$M_1 - M_0 = 1000$ cycles: $\delta = 100\,\mu\mathrm{s}$, far outside
any practical $m/z$ cell. The replayed signal is rejected by structure,
not by any filtering algorithm.

% =========================================================
\section{Discussion}
\label{sec:discussion}
% =========================================================

\subsection{Comparison with existing acquisition paradigms}

DDA interprets a full spectrum before selecting precursors. This
interpretation step is the root of both slowness and vulnerability:
slow because interpretation requires a full scan; vulnerable because
off-target ions that resemble targets can fool the intensity ranker.
Temporal programming eliminates the interpretation step entirely.
The trigger decision requires only $O(1)$ cell membership lookup.

SRM/PRM filters on $m/z$ after some form of mass selection. The
temporal programming equivalent is a pre-compiled cell registry,
but the key difference is the derivation: the cell registry is derived
from first principles via the four-analyzer compilers
(Section~\ref{sec:compiler}), giving an exact relationship between the
target window and the timing cell. No empirical window optimization
is required.

DIA/SWATH acquires complete fragment spectra for downstream
deconvolution. Temporal programming targets specific transitions;
it is complementary to DIA, not competitive. DIA provides comprehensive
coverage; temporal programming provides structural specificity for
known targets. In practice, the two paradigms can be combined:
a wide timing cell covers the full $m/z$ range (DIA mode) while a
narrow sub-registry covers specific targets (targeted mode),
with the trigger switching between modes by swapping $\Gamma$.

\subsection{The partition uncertainty as a design constraint}

The partition uncertainty relation~\eqref{eq:uncertainty} sets an
absolute limit on the combination of timing speed and mass resolution.
This limit is currently approached empirically through instrument
optimization; the framework derives it analytically. The practical
implication is that any claim of ``real-time'' triggering on a
sub-0.1\,Da window requires an acquisition time of at least
$\tau_{\min} \sim 100\,\mathrm{ms}$ (Table~\ref{tab:resolution}),
regardless of electronics speed or algorithm efficiency. This is a
fundamental constraint, not an engineering limitation.

\subsection{Implementation pathway}

The temporal programming acquisition system requires three hardware
components: a reference oscillator (already present in all
Orbitrap/FT-ICR instruments), an ISR that latches $\trec(k)$ at ion
arrival (requires $<100\,\mathrm{ns}$ interrupt latency, achievable
on current embedded processors), and a read-only cell registry (a
balanced interval tree in ROM, lookup in $O(\log m)$). No new
detector hardware is needed. The paradigm is a firmware/software
change to existing instruments.

% =========================================================
\section{Conclusion}
\label{sec:conclusion}
% =========================================================

We have established a unified first-principles framework for real-time
targeted acquisition in mass spectrometry grounded in the bounded phase
space formulation of ion dynamics and the temporal programming paradigm.

The central result is the \emph{Structural Incorruptibility Theorem}:
because the temporal programming MS system makes accept/discard decisions
by timing cell membership alone, with no $m/z$ parsing step, the system
has zero injection surface and off-target ions are structurally rejected
rather than statistically suppressed.

The \emph{partition uncertainty relation} $\Delta\Mop \cdot \tp \geq \hbar$,
derived from first principles via the partition Lagrangian, provides the
exact real-time/resolution design curve: $\tau_{\min} = \hbar/\delta\Mop_i$.
This curve is currently set empirically; it is now derivable.

The \emph{Composition-Inflation Theorem} $T(n,d) = d(d+1)^{n-1}$ shows
that adding a second timing axis (ion mobility drift time) produces a
state space of $2 \cdot 3^{n-1}$ categories, enabling isobar separation
that is structurally impossible on a single timing axis.

The four mass analyzer equations --- TOF, quadrupole, Orbitrap, FT-ICR ---
all emerge from a single partition Lagrangian, establishing that these
historically separate technologies are boundary conditions of the same
physical law. This unification provides the cell compilers that translate
$m/z$ target lists into timing registries, closing the loop from
physics to acquisition software.

\bibliographystyle{plainnat}
\bibliography{references}

\end{document}
